\section{实验概述}
本次实验围绕类 FAT-16 文件系统的设计与并发访问控制展开，主要包括两个任务。
任务一要求在内存中构造一个简化的文件系统，将数组 \texttt{Disk[DISK\_MAXLEN]} 模拟为硬盘空间，并在此基础上实现目录创建、目录查看、文件删除、文件打开、文件读取、文件写入和文件关闭等基本文件系统操作。
任务二则在任务一的基础上进一步考虑多个进程同时访问文件和目录时的同步问题，通过信号量机制保证文件系统在并发访问过程中的正确性。
\subsection{任务一}
在任务一中，实验的重点是完成文件系统基本结构的设计，包括磁盘空间划分、FAT 表维护、目录项组织以及文件数据块管理。通过将内存数组划分为超级块、FAT 表、根目录区和数据区，可以较为清晰地模拟真实文件系统中文件和目录的存储方式。
其中，FAT 表用于记录数据块之间的链接关系，目录项用于保存文件名、文件类型、起始块号和文件大小等元信息。通过这些结构，可以实现文件创建、写入、读取和删除等操作，并观察文件内容在模拟磁盘中的组织过程。

\subsection{任务二}
在任务二中，实验进一步引入多进程访问场景。由于多个进程可能同时访问同一个文件或目录，因此需要通过信号量对共享资源进行保护。
对于目录项、FAT 表和磁盘数据区等全局共享结构，需要使用互斥锁避免多个进程同时修改造成数据不一致；
对于普通文件的访问，则需要满足一个文件可以被多个进程同时读取，但同一时刻只能被一个进程写入的要求（类似于之前的读者-写者问题）。
因此，本任务的关键在于设计合理的读写同步机制，使文件系统既能保证并发访问的安全性，又能允许多个读进程并发执行，提高系统资源利用率。